\documentclass[10pt,a4paper]{amsart}
\usepackage[margin=19mm]{geometry}
\usepackage{amsmath,amssymb,mathtools}
\usepackage[scr=rsfs,cal=euler]{mathalfa}
\usepackage{xfrac}
\usepackage[unicode,hidelinks,bookmarks=false]{hyperref}
\newcommand{\ie}{i.e.\ }
\newcommand{\cf}{cf.\ }
\newcommand{\resp}{respectively\ }
\newcommand{\Subsection}{\subsection}
\providecommand{\nobreakdash}{\nobreak\hbox{-}}
\setlength{\emergencystretch}{2em}
\multlinegap0pt
\usepackage{amssymb}
\usepackage{dsfont}
\usepackage{enumerate}
\usepackage{xcolor}
\usepackage{esint}
\usepackage{mathtools}
\usepackage{tensor}
\usepackage[normalem]{ulem}
\usepackage{cancel}
\newtheorem{lemma}{Lemma}
\newcommand{\RR}{\mathds{R}}
\newcommand{\supp}{\mathrm{supp }}
\title{Regularity of Einstein $5$-manifolds via $4$-dimensional gap theorems}
\author{Yiqi Huang, Tristan Ozuch}
\date{Source: arXiv:2512.21317v2 (CC BY 4.0)}
\begin{document}
\maketitle

\noindent\emph{Source and licence.} Regularity of Einstein $5$-manifolds via $4$-dimensional gap theorems. Version arXiv:2512.21317v2 (CC BY 4.0), distributed by arXiv under the Creative Commons Attribution 4.0 International licence, http://creativecommons.org/licenses/by/4.0/, as recorded in the arXiv licence metadata for this exact version and shown on the abstract page https://arxiv.org/abs/2512.21317v2. Full licence notice: \texttt{licenses/arxiv-2512.21317-CC-BY-4.0.txt}.\par\smallskip

\noindent\textit{Editorial scope.} Counted unit: the article's lemma l:Caccioppoli (source0.tex lines 2005--2010). The article's complete original proof of it is delivered with it (source0.tex lines 2012--2076, one \texttt{proof} environment of 65 lines). No mathematics is written, repaired or re-derived: the statement and the proof are the article's own lines, in the article's own notation, and no retained line is substituted or re-flowed.  Retained uncounted as the source-local context the proof needs: lines 1997--1999.  Nothing was omitted from any retained span, so the package carries no omission record.  No retained line contains a citation, so the wrapper carries no bibliography and no bibliography entry.  No retained line contains a figure environment, an image-inclusion command or drawing code, so no image is delivered and the package declares no asset dependency.  Editorial formatting only, in addition: an AMS \texttt{amsart} preamble that restates the article's own declarations and macros used by the retained lines, namely the environments \texttt{lemma} on separate counters, together with the standard packages those lines need.  The retained environments print in this document as Lemma 1 in order of appearance, whereas the article prints the same items with the numbering of its own full document.  The complete native source of the article as distributed in the arXiv e-print for this exact version is bundled unchanged as \texttt{sources/191643/source0.tex}.
\begin{equation}\label{e: sub-equ of u}
    \Delta u \geqslant - f u. 
\end{equation}

\begin{lemma}\label{l:Caccioppoli}
    Let $u \in C^1(B^*)$ and $u \geqslant 0$ satisfying the equation \eqref{e: sub-equ of u} in the sense of distributions. Suppose $u \in L^{2q}$ for some $q \geqslant 1 + \varepsilon >1$ and $|f(x,y) | \leqslant \delta |y|^{-2}$ for some $\delta \leqslant \delta_0(\varepsilon,q)$. Then for any $\eta \in C^{\infty}_c(B)$ we have
    \begin{align}\label{e:Cacchio-1}
        \int_B \eta^2 |\nabla u^q|^2 \leqslant C(\varepsilon) q \int_B |\nabla \eta|^2 u^{2q}.
    \end{align}
\end{lemma}

\begin{proof}
    By integration by parts, the following holds for any function $\chi \in C_c^{\infty}(B^*)$ with compact support in $B^*$
    \begin{equation}\label{e:[Curv. Est] Int by part}
        \int_B \nabla \chi \cdot \nabla u \leqslant \int_B \chi f u.
    \end{equation}

We define the following test function for some large constant $K>0$
\begin{equation}
    F(t) = \frac{t^{q}}{(1 + \frac{t}{K})^{q/2}}
\end{equation}
whose motivation is that if $u\in L^{2q}$, then $F(u)\in L^4$ since we temper the large values of the function.

Note that the function $F$ is $C^{\infty}$ for $t>0$. We define $G(t) \equiv F(t) F'(t)$. Then the following properties hold by direct computation:
\begin{align}
    F(t) &\leqslant  K^{q/2} t^{q/2};\\
    tG(t) &\leqslant qF(t)^2;\label{e:[Curv. Est] tG < F}\\
    G'(t) &\geqslant c(\varepsilon) F'(t)^2. \label{e:[Curv. Est] G' > c F'}
\end{align}


Now we pick $\eta, \bar{\eta}_{\varepsilon} \in C_c^{\infty}(B)$ and further let $\bar{\eta}_{\varepsilon}$ vanish in $B_1^x \times B_{\varepsilon}^y$. Choose the cutoff function $\chi \equiv (\eta \bar{\eta}_{\varepsilon})^2 G(u)$ and by \eqref{e:[Curv. Est] Int by part} and \eqref{e:[Curv. Est] tG < F} we have
\begin{align}\label{eq:control nablaG nablau}
    \int_{B} \nabla \Big( (\eta \bar{\eta}_{\varepsilon})^2 G(u)\Big) \cdot \nabla u \leqslant \int_{B} (\eta \bar{\eta}_{\varepsilon})^2 G(u) f u \leqslant q \int_{B} (\eta \bar{\eta}_{\varepsilon})^2  |f| |F(u)|^2.
\end{align}

Note that $\nabla F(u) = F'(u) \nabla u$ and $\nabla G(u)= G'(u) \nabla u$. Hence by \eqref{e:[Curv. Est] G' > c F'} we have $\nabla G(u) \cdot \nabla u = |\nabla u|^2 G'(u) \geqslant c(\varepsilon) F'(u)^2 |\nabla u|^2 = c(\varepsilon)|\nabla F(u)|^2$. Hence
\begin{align}
    \int_{B} \nabla \Big( (\eta \bar{\eta}_{\varepsilon})^2 G(u)\Big) \cdot \nabla u &= \int_B (\eta \bar{\eta}_{\varepsilon})^2 \nabla G(u) \cdot \nabla u + \int_B G(u) \nabla (\eta \bar{\eta}_{\varepsilon})^2 \cdot \nabla u \\
    &\geqslant c(\varepsilon) \int_B (\eta \bar{\eta}_{\varepsilon})^2 |\nabla F(u)|^2- \int_B    |2 \eta \bar{\eta}_{\varepsilon} F'(u) \nabla u| \cdot  |  F(u) \nabla(\eta \bar{\eta}_{\varepsilon}) |   \\
    &= c(\varepsilon) \int_B (\eta \bar{\eta}_{\varepsilon})^2 |\nabla F(u)|^2- \int_B    |2 \eta \bar{\eta}_{\varepsilon} \nabla F(u)| \cdot  |  F(u) \nabla(\eta \bar{\eta}_{\varepsilon}) |
\end{align}


Combining these and \eqref{eq:control nablaG nablau} we obtain
\begin{align}\label{e:[Curv. Est] I_1 + I_2 inequality}
    c(\varepsilon) \int_B (\eta \bar{\eta}_{\varepsilon})^2 |\nabla F(u)|^2 \leqslant \int_{B} |2 \eta \bar{\eta}_{\varepsilon} \nabla F(u)| \cdot |\nabla(\eta \bar{\eta}_{\varepsilon}) F(u)|  + q \int_B (\eta \bar{\eta}_{\varepsilon})^2 |f| F(u)^2 \equiv I_1 + I_2.
\end{align}

By Cauchy-Schwarz inequality, for some $c'(\varepsilon)$ we have
\begin{equation}
    I_1 \leqslant \frac{c(\varepsilon)}{2} \int_B (\eta \bar{\eta}_{\varepsilon})^2 |\nabla F(u)|^2 + c'(\varepsilon) \int_B |\nabla (\eta \bar{\eta}_{\varepsilon})|^2 F(u)^2.
\end{equation}

To estimate $I_2$, we apply Hardy inequality to get
\begin{equation}\label{e:[curvatuer bound] I_2}
    I_2 \leqslant  \delta q \int_{B_1^x} \int_{B_1^y} 
    (\eta \bar{\eta}_{\varepsilon})^2 F(u)^2 |y|^{-2} dy dx \leqslant  4\delta q \int_{B_1^x} \int_{B_1^y} |\nabla_y   (\eta \bar{\eta}_{\varepsilon} \cdot F(u))|^2(x,y) dy dx. 
\end{equation}

Note that $|\nabla_y F(u)|^2 \leqslant |\nabla F(u)|^2$. Hence by choosing $\delta \leqslant \delta_0(\varepsilon,q)$ small, we can absorb all terms involving $|\nabla F|$ into the left hand side of \eqref{e:[Curv. Est] I_1 + I_2 inequality} and thus obtain
\begin{equation}
    \int_{B} (\eta \bar{\eta}_{\varepsilon})^2 |\nabla F(u)|^2 \leqslant C(\varepsilon)q \int_B |\nabla (\eta \bar{\eta}_{\varepsilon})|^2 F^2(u).
\end{equation}
Here is the only place where we require $\delta$ to be small depending on $q$. 

Note that we can choose $\bar{\eta}_{\varepsilon}$ such that $\eta \cdot \bar{\eta}_{\varepsilon} \to \eta$ and $\int_{B \cap \supp(\eta)} | \nabla \bar{\eta}_{\varepsilon}|^4 \to 0$ as $\varepsilon \to 0$ since $B\setminus \supp(\Bar{\eta}_{\varepsilon}) \subset B_1^x \times B_{2\varepsilon}^y \to B_1^x \times \{0\} \subset \RR^{n-4} \times \{0\}$. Combining $F(u) \leqslant K^{q/2} u^{q/2}$ and $u \in L^{2q}$, we have $F(u) \in L^4$ and thus
\begin{equation}
    \int_{B} |\nabla \bar{\eta}_{\varepsilon}|^2 \eta^2 F^2(u) \leqslant C(K) \Bigg(\int_{B} |\nabla \bar{\eta}_{\varepsilon}|^4 \Bigg)^{1/2} \cdot \Bigg( \int_B F(u)^4 \Bigg)^{1/2} \to 0, \text{ as } \varepsilon \to 0.
\end{equation}

Combining all, we can now conclude that
\begin{equation}
    \int_{B} \eta^2 |\nabla F(u)|^2 \leqslant C(\varepsilon)q \int_B |\nabla \eta|^2 F^2(u).
\end{equation}
Finally we let $K \to \infty$. Note that $F(u) \to u^{q}$. The proof is finished by Fatou's Lemma.    \end{proof}

\end{document}
